\Askhsh{28338_SOLUTION}{1}
\begin{ylh}{1}
    \item [Ύλη από εικόνα]
\end{ylh}
ΛΥΣΗ
\begin{alist}
    \item \begin{rlist}
        \item Η συνάρτηση $f$ είναι παραγωγίσιμη στο $\mathbb{R}$ και παρουσιάζει τοπικό ελάχιστο στο $x_1=2$, οπότε από το θεώρημα του Fermat θα είναι $f'(2)=0$. Δίνεται επίσης ότι $f''(2)=-32$. Επομένως, η εξίσωση της εφαπτομένης της $C_f$ στο σημείο με τετμημένη $x_1=2$, είναι η
        \[ y - f(2) = f'(2)(x - 2) \]
        \[ y + 32 = 0(x - 2) \]
        \[ y = -32. \]
        \item Οι γραφικές παραστάσεις της $f$ και της $f'$ τέμνονται στο σημείο $A(-2,0)$, άρα $f(-2)=0$ και $f'(-2)=0$. Επομένως, η εξίσωση της εφαπτομένης της $C_{f'}$ στο σημείο με τετμημένη $x_2=-2$, είναι η
        \[ y - f'(-2) = f''(-2)(x + 2) \]
        \[ y - 0 = 0(x + 2) \]
        \[ y = 0. \]
    \end{rlist}
    \item \begin{rlist}
        \item Στο (α) ερώτημα βρήκαμε $f'(2)=f'(-2)=0$ και επειδή η $f'$ είναι πολυωνυμική συνάρτηση $2^{ου}$ βαθμού, θα έχει ακριβώς δύο ρίζες, τις $x_1=2$ και $x_2=-2$. Επομένως είναι
        \[ f'(x) = \alpha(x-x_1)(x-x_2) = \alpha(x-2)(x+2) = \alpha(x^2-4) = \alpha x^2 - 4\alpha, \quad \text{με } \alpha \ne 0. \]
        Η γραφική παράσταση της $f'$ διέρχεται από το σημείο $B(0,-12)$, άρα είναι $f'(0)=-12$, οπότε $\alpha \cdot 0^2 - 4\alpha = -12$ ή $\alpha=3$.

        Επομένως είναι
        \[ f'(x) = 3x^2 - 12, \quad x \in \mathbb{R}. \]
        \item Για κάθε $x \in \mathbb{R}$ είναι $f'(x) = 3x^2 - 12 = (x^3 - 12x)'$, άρα $f(x) = x^3 - 12x + c$, $c \in \mathbb{R}$.

        Όμως
        \[ f(2) = -32 \]
        ή $8 - 24 + c = -32$
        ή $c = -16$.

        Επομένως είναι
        \[ f(x) = x^3 - 12x - 16, \quad x \in \mathbb{R}. \]
        \item Οι ρίζες της $f'$ είναι οι $x_1 = -2$, $x_2 = 2$ και το πρόσημο της $f'$ φαίνεται στον παρακάτω πίνακα, από τον οποίο προσδιορίζουμε τα διαστήματα μονοτονίας και τα τοπικά ακρότατα της $f$. Επίσης είναι
        \[ \lim_{x \to -\infty} f(x) = \lim_{x \to -\infty} (x^3 - 12x - 16) = \lim_{x \to -\infty} x^3 = +\infty \]
        και
        \[ \lim_{x \to +\infty} f(x) = \lim_{x \to +\infty} (x^3 - 12x - 16) = \lim_{x \to +\infty} x^3 = +\infty. \]
    \end{rlist}
\end{alist}
\newpage
\begin{tabular}{|c||c|c|c|c|c|c|c|}
\hline
$x$ & $-\infty$ & & $-2$ & & $2$ & & $+\infty$ \\
\hline
$f'$ & \multicolumn{1}{|c|}{} & $+$ & $0$ & $-$ & $0$ & $+$ & \multicolumn{1}{|c|}{} \\
\hline
$f$ & $+\infty$ & $\uparrow$ & $32$ & $\downarrow$ & $-32$ & $\uparrow$ & $+\infty$ \\
\cline{3-3}\cline{5-5}
\multicolumn{1}{|c||}{} & & & \text{T.M.} & & \text{T.E.} & & \\
\hline
\end{tabular}

Η $f$ είναι συνεχής στα διαστήματα $\Delta_1 = (-\infty, -2]$, $\Delta_2 = [-2, 2]$, $\Delta_3 = [2, +\infty)$ και επειδή είναι γνησίως αύξουσα στα $\Delta_1$, $\Delta_3$ και γνησίως φθίνουσα στο $\Delta_2$, θα είναι
\[ f(\Delta_1) = \left(\lim_{x \to -\infty} f(x), f(-2)\right] = (-\infty, 32], \quad f(\Delta_2) = [f(2), f(-2)] = [-32, 32] \]
και $f(\Delta_3) = [f(2), \lim_{x \to +\infty} f(x)) = [-32, +\infty)$.
Το $-20$ ανήκει στα διαστήματα $f(\Delta_1)$, $f(\Delta_2)$ και $f(\Delta_3)$, άρα η εξίσωση $f(x)=-20$ έχει μία τουλάχιστον ρίζα σε καθένα από τα διαστήματα $\Delta_1$, $\Delta_2$ και $\Delta_3$.

Επιπλέον η $f$ είναι γνησίως μονότονη στα $\Delta_1$, $\Delta_2$ και $\Delta_3$, άρα υπάρχει ακριβώς μία ρίζα σε καθένα από τα διαστήματα $\Delta_1$, $\Delta_2$ και $\Delta_3$.

Άρα, η εξίσωση $f(x) = -20$ έχει ακριβώς τρεις διαφορετικές πραγματικές ρίζες.

\end{rlist}